\subsection{A fixed countable target for unimodular extensions}

\begin{claim}{A.18}{Proved}{Unimodular towers retain a nontrivial countable image}
Let $H$ be a nontrivial countable torsion-free group. There are a
countable torsion-free group $U$ and an embedding $j:H\hookrightarrow U$
with the following properties.
\begin{enumerate}
\item The inclusion induces isomorphisms
$H_k(H;\mathbb Z)\xrightarrow{\cong}H_k(U;\mathbb Z)$ for all $k\geq0$.
Every equation $w(t)=1$ with $w\in U*\langle t\rangle$ and total
exponent of $t$ equal to $1$ or $-1$ has a solution in $U$.
\item For every CW model $K=BH$, there is a model $L=BU$ containing
$K$ as a subcomplex and obtained from it by countably many pairs of
one- and two-cells. In particular,
$\dim L\leq\max(\dim K,2)$. If $K$ is countable then $L$ is countable.
\item Let $\tau$ be any small ordinal and let $(G_\alpha)_{\alpha\leq\tau}$
be a continuous tower beginning with $G_0=H$. At each successor suppose
\[
 G_{\alpha+1}
   =\bigl(G_\alpha*\langle t_\alpha\rangle\bigr)/
      \langle\!\langle w_\alpha(t_\alpha)\rangle\!\rangle,
 \qquad
 \operatorname{exp}_{t_\alpha}(w_\alpha)=\pm1,
 \tag{A.18.1}
\]
with the canonical structure map, and at each limit take the group
colimit. Then there is a homomorphism
$\Phi:G_\tau\to U$ with $\Phi|_H=j$. Its image is a nontrivial
countable quotient of $G_\tau$.
\end{enumerate}
Consequently a transfinite construction consisting of these extensions,
starting from a nontrivial countable group with a two-dimensional
classifying space, cannot have the no-countable-image property in
A3(iii). The conclusion covers finite layered systems after subdividing
them into their successive generator--relator pairs.
\end{claim}

\begin{proof}
We first check the geometric input and the specific maps it supplies.
For a torsion-free group $P$ and a word $w\in P*\langle t\rangle$
of total exponent $\epsilon=\pm1$, put
\[
 P_w=(P*\langle t\rangle)/\langle\!\langle w\rangle\!\rangle.
\]
Choose a classifying CW complex $B$ for $P$, attach one based loop
representing $t$, and then one two-cell along a loop representing $w$.
Write the resulting complex as $B_w$. Forester--Rourke
\cite[Theorem~1 and Proposition~1]{CpathsForesterRourke2005Adjunction}
prove that the specified homomorphism $P\to P_w$ is injective,
$P_w$ is torsion-free, and the inclusion-induced map
\[
 \mathbb Z[P_w]\otimes_{\mathbb Z[P]}\pi_2(B)
       \longrightarrow \pi_2(B_w)
 \tag{A.18.2}
\]
is an isomorphism. Their one-pair result applies here because the
one-by-one exponent matrix is $(\epsilon)$. The diagram proof
\cite[\S\S1--3]{CpathsForesterRourkeDiagrams}
first proves surjectivity of
$\pi_2(B\cup e^1,B)\to\pi_2(B_w,B)$ by removing cancelling
diagram pairs; its Lemma~3.1 and Proposition~3.3 give precisely
the injection and the map in (A.18.2).

Here is why this gives the classifying space needed in our construction.
The preimage of $B$ in the universal cover of $B_w$ consists of copies
of the contractible universal cover of $B$, since $P\to P_w$ is
injective. Collapse each such component to a vertex. These collapses
are a homotopy equivalence, and the resulting space is a
two-dimensional CW complex: outside those components only the new
one- and two-cells occur. It is simply connected, and (A.18.2)
makes its second homotopy group zero. Hurewicz and Whitehead
therefore make it contractible. Thus $B_w$ is a classifying
space for $P_w$. The actual relative cellular chain complex
with constant integral coefficients is
\[
 C_*(B_w,B;\mathbb Z)
       =\bigl(0\longrightarrow\mathbb Z
           \xrightarrow{\epsilon}\mathbb Z\longrightarrow0\bigr)
\]
in degrees two and one. It is acyclic. The inclusion $B\subset B_w$
therefore induces the stated ordinary integral homology isomorphisms
and retains dimension at most $\max(\dim B,2)$.

Now set $U_0=H$. Suppose the countable torsion-free group $U_n$
has been constructed. List all finite words in
$U_n*\langle t\rangle$ having total $t$-exponent $\pm1$ as
$w_{n,0},w_{n,1},\ldots$. There are only countably many. Starting
with $U_{n,0}=U_n$, form
\[
 U_{n,k+1}
  =\bigl(U_{n,k}*\langle t_{n,k}\rangle\bigr)/
     \langle\!\langle w_{n,k}(t_{n,k})\rangle\!\rangle,
 \qquad
 U_{n+1}=\mathop{\rm colim}_{k<\omega}U_{n,k}.
 \tag{A.18.3}
\]
The coefficients of $w_{n,k}$ lie in the specified subgroup $U_n$
of $U_{n,k}$. Every step is consequently covered by the preceding
paragraph. All structure maps are injective, all groups are
torsion-free, and each remains countable. Set
$U=\mathop{\rm colim}_{n<\omega}U_n$ and let $j$ be the composite
inclusion of $H$.

Perform (A.18.3) on CW complexes by the actual one- and two-cell
attachments just described, and take the nested unions at the
two countable limit stages. Every map from a sphere of dimension at least
two into one of
these unions has image in a finite subcomplex, hence factors
through an earlier stage. Since that stage is aspherical, the
map is nullhomotopic. The union is therefore aspherical and
is a model for the indicated group colimit. Only countably
many cells have been added to $K$, in dimensions one and two.
The ordinary cellular homology of a nested CW union is its
filtered homology colimit, so the homology isomorphisms persist
from $H$ to $U$. This proves the first homological assertion and
all dimensional and cell-count assertions in part~2.

For the equation assertion, take a word $w\in U*\langle t\rangle$
with exponent $\pm1$. Its finitely many coefficients lie in one
$U_n$. The word therefore occurs as some $w_{n,k}$ in the chosen
list. The image of $t_{n,k}$ in $U_{n+1}\subset U$ is an actual
solution. Thus all the equations asserted in part~1 have solutions
inside this one fixed group $U$.

It remains to construct the map of part~3, including its limit
compatibilities. We use transfinite recursion. Put $\phi_0=j$.
Given $\phi_\alpha:G_\alpha\to U$, apply it to all coefficients
of $w_\alpha$. The resulting word
\[
 \overline w_\alpha(t)
       =(\phi_\alpha*{\rm id}_{\langle t\rangle})(w_\alpha(t))
       \in U*\langle t\rangle
\]
still has total exponent $\pm1$, even if some coefficients become
trivial and the word reduces. Choose a solution $u_\alpha\in U$.
The homomorphism on the free product given by $\phi_\alpha$ on
$G_\alpha$ and $t_\alpha\mapsto u_\alpha$ kills the specified
relator, hence descends to
$\phi_{\alpha+1}:G_{\alpha+1}\to U$. It satisfies
$\phi_{\alpha+1}|_{G_\alpha}=\phi_\alpha$ exactly. At a limit ordinal
$\beta$, the group-colimit universal property gives the unique
$\phi_\beta:G_\beta\to U$ with these restrictions.
This constructs $\Phi=\phi_\tau$. Since $\Phi|_H=j$ is injective
and $H\ne1$, its image is nontrivial; it is countable because it
is a subgroup of $U$. Factoring $\Phi$ through its image gives
the claimed explicit countable quotient.

For a finite layered system the exponent matrix is triangular
in its specified order. Integral invertibility makes each
diagonal entry $\pm1$, so each individual layer is of the form
(A.18.1). This proves the last assertion for that class as well.
The proof uses the specified layering and does not supply a
map to $U$ for an arbitrary simultaneous system whose exponent
matrix is invertible.
\end{proof}

